% Generated draft: requires mathematical and source review.
\documentclass[11pt]{article}
\usepackage[margin=25mm]{geometry}
\usepackage{fontspec}
\setmainfont{Latin Modern Roman}
\usepackage{amsmath,amssymb,mathtools,unicode-math}
\setmathfont{Latin Modern Math}
\usepackage{enumitem,needspace,xurl,hyperref}
\setlength{\parindent}{0pt}
\setlength{\parskip}{5pt}
\setlength{\emergencystretch}{3em}
\newcommand{\N}{\mathbb{N}}
\newcommand{\Z}{\mathbb{Z}}
\newcommand{\Q}{\mathbb{Q}}
\newcommand{\R}{\mathbb{R}}
\newcommand{\C}{\mathbb{C}}
\newcommand{\F}{\mathbb{F}}
\newcommand{\A}{\mathbb{A}}
\newcommand{\PP}{\mathbb{P}}
\newcommand{\E}{\mathbb{E}}
\newcommand{\eps}{\varepsilon}
\newcommand{\dd}{\,\mathrm{d}}
\begin{document}
\section*{UnsolvedMath ID: 30004398}
\textbf{Source problem code:} catalogue entry
\subsection{Definitions and mathematical statement}
The supplied question is: ``Is the Bestvina complex an equivariant deformation retract of the Davis complex?'' The record does not specify the acting group, the constructions of the two complexes, their hypotheses, or an inclusion between them. Consequently, the following gives the precise retraction condition once that missing context is supplied; it does not select or assert an intended class of examples.

Let \(G\) be a group acting continuously on spaces \(\Sigma\) and \(B\), intended to denote the Davis complex and the Bestvina complex in the unspecified source setting. Suppose a \(G\)-equivariant topological embedding \(i\colon B\hookrightarrow\Sigma\) has been specified. To say that \(i(B)\) is a \(G\)-equivariant deformation retract of \(\Sigma\) means that there exist continuous maps
\[
r\colon\Sigma\longrightarrow B,
\qquad H\colon\Sigma\times[0,1]\longrightarrow\Sigma
\]
such that
\[
\begin{aligned}
r(i(b))&=b &&(b\in B),\\
r(gx)&=g\,r(x) &&(g\in G,\ x\in\Sigma),\\
H(x,0)&=x &&(x\in\Sigma),\\
H(x,1)&=i(r(x)) &&(x\in\Sigma),\\
H(i(b),t)&=i(b) &&(b\in B,\ t\in[0,1]),\\
H(gx,t)&=g\,H(x,t) &&(g\in G,\ x\in\Sigma,\ t\in[0,1]).
\end{aligned}
\]
Here the action on \([0,1]\) is trivial. Thus the conditional question is whether such \(r\) and \(H\) exist for the intended complexes and inclusion. This uses the convention that a deformation retraction fixes the retract throughout the homotopy; the supplied record does not establish the source's convention or specify whether an inclusion is prescribed or is to be constructed.

The registry marks this entry as \texttt{unrecoverable} and identifies it as a duplicate of the second question represented canonically by \texttt{OWR-17470-013}. It explicitly promotes no separate clean statement. The displayed formulation therefore remains conditional and requires recovery of the source context.
\subsection{Short English statement}
The question asks whether the Davis complex can be continuously deformed onto the Bestvina complex, compatibly with the group action and while fixing the retract throughout. The acting group, constructions, hypotheses, and inclusion are unspecified. This retained duplicate has no separately verified clean statement; the supplied open status and literature assessment require independent verification.
\subsection{Sources}
\begin{itemize}
\item Local UnsolvedMath registry, permanent ID \(30004398\), exact source code \texttt{catalogue entry}, title ``Duplicate Bestvina-Complex Retraction Question.'' The registry retains the original question for provenance, marks its statement status as \texttt{unrecoverable}, and identifies it as repeating the second question represented canonically by \texttt{OWR-17470-013}. No proposer is supplied; the proposed year is recorded as 2020.
\item Supplied source citation: \emph{Manifolds and Groups (2021)}, Oberwolfach Reports, DOI 10.4171/owr/2020/7, pp.\ 417--463. \url{https://doi.org/10.4171/owr/2020/7}. The background describes these reports as documenting mathematical workshops and their open-problem discussions. The supplied citation's 2021 label, DOI containing 2020, and recorded proposed year 2020 are retained without reconciliation.
\item Supplied literature reference: Petrosyan--Prytula 2020. \url{https://doi.org/10.4171/GGD/581}. No further bibliographic details are supplied.
\item Literature assessment dated exactly 2026-08-22: the supplied curation labels the problem open and reports equivariant deformation retractions under additional subconical or low-dimensional hypotheses in Petrosyan--Prytula, while describing the unrestricted Davis-to-Bestvina deformation-retract question as unresolved. The listed literature sources are Petrosyan--Prytula 2020, \url{https://doi.org/10.4171/GGD/581}, and the Oberwolfach report, \url{https://doi.org/10.4171/owr/2020/7}. This is supplied dated curation derived from the OWR record, not an independently verified literature assessment; both the time-sensitive status and cited claims require independent verification.
\end{itemize}
\end{document}
